% Options for packages loaded elsewhere
\PassOptionsToPackage{unicode}{hyperref}
\PassOptionsToPackage{hyphens}{url}
%
\documentclass[
  11pt,
]{article}
\usepackage{amsmath,amssymb}
\usepackage{iftex}
\ifPDFTeX
  \usepackage[T1]{fontenc}
  \usepackage[utf8]{inputenc}
  \usepackage{textcomp} % provide euro and other symbols
\else % if luatex or xetex
  \usepackage{unicode-math} % this also loads fontspec
  \defaultfontfeatures{Scale=MatchLowercase}
  \defaultfontfeatures[\rmfamily]{Ligatures=TeX,Scale=1}
\fi
\usepackage{lmodern}
\ifPDFTeX\else
  % xetex/luatex font selection
\fi
% Use upquote if available, for straight quotes in verbatim environments
\IfFileExists{upquote.sty}{\usepackage{upquote}}{}
\IfFileExists{microtype.sty}{% use microtype if available
  \usepackage[]{microtype}
  \UseMicrotypeSet[protrusion]{basicmath} % disable protrusion for tt fonts
}{}
\makeatletter
\@ifundefined{KOMAClassName}{% if non-KOMA class
  \IfFileExists{parskip.sty}{%
    \usepackage{parskip}
  }{% else
    \setlength{\parindent}{0pt}
    \setlength{\parskip}{6pt plus 2pt minus 1pt}}
}{% if KOMA class
  \KOMAoptions{parskip=half}}
\makeatother
\usepackage{xcolor}
\usepackage[margin=1in]{geometry}
\usepackage{longtable,booktabs,array}
\usepackage{calc} % for calculating minipage widths
% Correct order of tables after \paragraph or \subparagraph
\usepackage{etoolbox}
\makeatletter
\patchcmd\longtable{\par}{\if@noskipsec\mbox{}\fi\par}{}{}
\makeatother
% Allow footnotes in longtable head/foot
\IfFileExists{footnotehyper.sty}{\usepackage{footnotehyper}}{\usepackage{footnote}}
\makesavenoteenv{longtable}
\setlength{\emergencystretch}{3em} % prevent overfull lines
\providecommand{\tightlist}{%
  \setlength{\itemsep}{0pt}\setlength{\parskip}{0pt}}
\setcounter{secnumdepth}{-\maxdimen} % remove section numbering
\ifLuaTeX
  \usepackage{selnolig}  % disable illegal ligatures
\fi
\usepackage{bookmark}
\IfFileExists{xurl.sty}{\usepackage{xurl}}{} % add URL line breaks if available
\urlstyle{same}
\hypersetup{
  pdftitle={Certified spectral distance for finite Jacobi matrices},
  hidelinks,
  pdfcreator={LaTeX via pandoc}}

\title{Certified spectral distance for finite Jacobi matrices}
\author{}
\date{}

\begin{document}
\maketitle

\section{1. Introduction}\label{introduction}

Consider the symmetric tridiagonal matrix

\[
J=\begin{pmatrix}0&1&0\\1&0&1\\0&1&0\end{pmatrix},
\qquad \tau=\frac75.
\]

The polynomial identity \(\det(\tau I-J)\ne0\) already decides whether
\(J-\tau I\) is invertible. It does not quantify how close the matrix is
to singularity. For solving a shifted linear system, that quantitative
question is often the consequential one: how strongly can the inverse
amplify an input error or a residual, and how much perturbation can the
conclusion tolerate?

A finite exact-rational computation answers it here. The characteristic
continuants show that \((139/100,141/100]\) contains no eigenvalue,
while \((141/100,71/50]\) contains at least one. Those two finite
counting facts imply

\[
\boxed{\frac1{100}\le
\operatorname{dist}\!\left(\frac75,\sigma(J)\right)
\le\frac1{50}.}
\tag{1.1}
\]

The first inequality proves that \(J-\tau I\) is invertible and that its
Euclidean inverse norm is at most \(100\). If \(x\) solves
\((J-\tau I)x=f\) and a proposed \(\widehat x\) has a separately
certified residual \(\|f-(J-\tau I)\widehat x\|_2\le\delta\), then
\(\|x-\widehat x\|_2\le100\delta\). Thus the certificate is not
principally an eigenvalue-finding device: it is finite evidence for a
\emph{quantitative safety margin}.

For an arbitrary finite real symmetric matrix, put

\[
\Delta(J,\tau)=\operatorname{dist}(\tau,\sigma(J))
=\min_{\lambda\in\sigma(J)}|\lambda-\tau|.
\tag{1.2}
\]

Whenever \(\Delta(J,\tau)>0\), the spectral theorem gives

\[
\|(J-\tau I)^{-1}\|_2=\Delta(J,\tau)^{-1}.
\tag{1.3}
\]

Our central problem is therefore to furnish rational numbers
\(0<L\le U\) and evidence establishing \(L\le\Delta(J,\tau)\le U\), with
the additional property that a \textbf{separate verifier}, using the
same exact input and the requested tolerance \(\varepsilon>0\), can
check \(U-L\le\varepsilon\). This last comparison cannot be replaced by
a producer's assurance that it ran enough iterations. Likewise, a proof
of termination is distinct from a demonstration that the implementation
is economical at numerically relevant tolerances.

We study unreduced finite Jacobi matrices with \emph{rational}
coefficients and a rational target energy. The natural matrix is
represented by its \(n\) diagonal entries and \(n-1\) nonzero
off-diagonal entries; any extension to infinite indexing in the proof
development is an implementation convention, not new input data.
Rationality makes every continuant sign and Sturm-count comparison
decidable without floating-point assumptions. It is not a claim that
general real coefficients or irrational phase energies are executable
inputs to the same exact-data producer.

The key construction brackets the \textbf{distance from the target},
rather than identifying a nearest eigenvalue. An inner symmetric
interval with equal Sturm counts at its endpoints certifies spectral
absence; a larger symmetric interval with a strict count drop certifies
presence. Count-preserving rational bisection halves the width of the
resulting distance enclosure. The inner and outer radii are computed
directly from continuant bounds and a finite coefficient budget. The
mathematical certificate is checked independently of that construction.

The ingredients of this scheme have substantial prior art. Sturm counts
and guaranteed tridiagonal eigenvalue enclosures are classical;
Fernando's exact-inertia algorithm gives rigorous eigenvalue inclusions
{[}1{]}. Untrusted production followed by formally verified checking is
likewise established in real-algebraic computation {[}2{]}, and Lean's
\texttt{hex-real-roots} implements certified half-open Sturm root
isolation with stronger general polynomial functionality {[}3{]}. We
neither claim a first Sturm certificate nor propose a general
replacement for such root isolators. The specific question addressed
here is how a directly checkable \emph{target-centered Jacobi witness}
can be constructed from the three-term recurrence and consumed as a
quantitative spectral-distance, inverse-norm, and stability margin.
Section 7 describes this distinction and its limitations in more detail.

The exposition separates a mathematical assertion from evidence that a
particular implementation meets it. Sections 2 and 3 state the exact
counting convention and the certificate's semantics. Section 4 proves
constructive initialization and finite refinement. Section 5 deduces
standard perturbation and residual bounds with explicitly represented
uncertainties. Section 6 examines arithmetic resonance in an exactly
solvable chain; it explains the demand for narrow margins without
asserting a universal algorithmic lower bound. Finally, Section 7
compares the certificate with existing inclusion and root-isolation
methods, and records the limits of its computational trust boundary.

\section{2. Jacobi data and exact spectral
counting}\label{jacobi-data-and-exact-spectral-counting}

\subsection{2.1. The finite unreduced
matrix}\label{the-finite-unreduced-matrix}

Fix an integer \(n\ge1\), rational numbers

\[
d_0,\ldots,d_{n-1}\in\mathbb Q,
\qquad b_0,\ldots,b_{n-2}\in\mathbb Q\setminus\{0\},
\tag{2.1}
\]

and define the real symmetric Jacobi matrix
\(J_n=J_n(d,b)\in\mathbb R^{n\times n}\) by

\[
(J_n)_{ij}=
\begin{cases}
 d_i,&i=j,\\
 b_i,&j=i+1,\\
 b_j,&i=j+1,\\
 0,&|i-j|\ge2.
\end{cases}
\tag{2.2}
\]

The condition \(b_i\ne0\) makes the matrix \emph{unreduced}. No sign
restriction is imposed on \(b_i\); in the continuant recurrence only
\(b_i^2\) occurs. For \(n=1\) there are no off-diagonal conditions. We
write the real eigenvalues of \(J_n\), repeated according to
multiplicity, as \(\lambda_1,\ldots,\lambda_n\). In the unreduced Jacobi
case they are in fact simple, but simplicity is not used in the
elementary deduction from a certified exclusion interval to a distance
bound.

All rational data are henceforth regarded as real numbers when used in
matrix or spectral expressions. This embedding is exact. In particular,
a theorem about \(\mathbb Q\)-valued computation need not assume that an
arbitrary real eigenvalue has a rational representation.

The formal recurrence is parameterized by functions
\(d,b:\mathbb N\to\mathbb Q\). For an \(n\)-site matrix, only the finite
prefixes in (2.1) are consulted. A practical serialized input may
therefore be represented by two finite arrays without altering the
mathematical object; that serialization itself is not part of the
spectral theorem proved here.

We consider a target \(\tau\in\mathbb Q\). The exact singularity
question is whether \(\tau\in\sigma(J_n)\); in the nonsingular case the
object to enclose is

\[
\Delta_n(\tau):=\min_{1\le j\le n}|\lambda_j-\tau|>0.
\tag{2.3}
\]

\subsection{2.2. Characteristic
continuants}\label{characteristic-continuants}

Let \(J_k\) denote the leading \(k\times k\) principal submatrix of
\(J_n\), with \(J_0\) the empty matrix. Set

\[
q_0(x)=1,\qquad q_1(x)=x-d_0,
\tag{2.4}
\]

and, for \(0\le k\le n-2\), define

\[
q_{k+2}(x)
=(x-d_{k+1})q_{k+1}(x)-b_k^{\,2}q_k(x).
\tag{2.5}
\]

Expanding a tridiagonal determinant along its final row or column gives

\[
q_k(x)=\det(xI_k-J_k),\qquad 0\le k\le n.
\tag{2.6}
\]

In particular, \(q_k\in\mathbb Q[x]\) is monic of degree \(k\). The
target \(\tau\) belongs to \(\sigma(J_n)\) exactly when \(q_n(\tau)=0\).
This exact rational equality is decidable; it is not a test involving a
numerical singular-value threshold.

For \(x\in\mathbb Q\), each value \(q_k(x)\) is rational and can be
computed from (2.4)--(2.5) using additions and multiplications alone.
There is no division by \(q_k(x)\). This matters when an intermediate
continuant vanishes: the recurrence remains defined, and its zero is
handled by the sign-count convention below rather than by a perturbation
of the test energy.

\subsection{2.3. A strict spectral-count
convention}\label{a-strict-spectral-count-convention}

For a finite list of real numbers, let \(\operatorname{var}\) be the
number of changes of sign after deleting all zero entries. At a rational
test energy \(x\), define

\[
N_n(x):=
\operatorname{var}\bigl(q_0(x),q_1(x),\ldots,q_n(x)\bigr).
\tag{2.7}
\]

Thus the list \((1,0,-2,3,3)\) has two variations. The deletion of zero
terms is part of the definition, not an informal convention that may be
altered by an implementation.

\textbf{Theorem 2.1 (finite Jacobi Sturm count).} \emph{For the
unreduced Jacobi matrix (2.2) and every real \(x\), the sign variation
count of the real continuants satisfies}

\[
N_n(x)=\#\{j\in\{1,\ldots,n\}:\lambda_j>x\}.
\tag{2.8}
\]

\emph{In particular, for rational \(a\le b\),}

\[
N_n(a)-N_n(b)
=\#\{j:a<\lambda_j\le b\}.
\tag{2.9}
\]

Here (2.8) denotes the corresponding real-variable sign count; (2.7) is
its exact rational restriction. The equality is the finite Sturm
oscillation theorem in the strict-greater-than convention used
throughout this paper.

\emph{Proof.} The Sturm oscillation theorem for an unreduced Jacobi
chain identifies the number of sign changes in its
leading-principal-minor continuant sequence, with zero values discarded,
with the number of eigenvalues strictly above the test energy. Applying
this fact at \(a\) and \(b\), and subtracting the two finite
cardinalities, leaves precisely the eigenvalues satisfying
\(a<\lambda_j\le b\). \(\square\)

Two consequences deserve emphasis. First, \(N_n\) is nonincreasing: if
\(a\le b\), then \(N_n(b)\le N_n(a)\). Second, equality of the counts at
two ordered endpoints means that the \emph{half-open} interval \((a,b]\)
contains no eigenvalue. A zero of \(q_n\) at \(b\) is included in that
interval, whereas a zero at \(a\) is excluded. These precise endpoints
will be retained in the certificate theorem; replacing \((a,b]\) by
\([a,b]\) or \((a,b)\) without adjusting the conclusion would change its
meaning.

\textbf{Proposition 2.2 (rational evaluation fidelity).} \emph{Evaluate
(2.4)--(2.5) over \(\mathbb Q\) at a rational \(x\), delete the zero
entries, and count their sign variations. The result agrees with the
sign variation count obtained by evaluating the corresponding real
characteristic polynomials at the image of \(x\) in \(\mathbb R\).}

\emph{Proof.} The embeddings \(\mathbb Q\hookrightarrow\mathbb R\)
preserve addition, multiplication, order, and the distinction between
zero and nonzero. A two-step induction on (2.5) therefore shows equality
of every rational continuant after casting with its real-polynomial
counterpart. The sign list is preserved entry by entry, and hence so is
its variation count. \(\square\)

Proposition 2.2 is the essential computational bridge: a checker can
work entirely with rationals and still invoke Theorem 2.1 about the
spectrum of a real matrix. It has no need to approximate an eigenvalue
or to compute a floating-point determinant.

\section{3. Certificates for spectral
distance}\label{certificates-for-spectral-distance}

\subsection{3.1. Exclusion and occupation as finite
evidence}\label{exclusion-and-occupation-as-finite-evidence}

A \emph{spectral-distance witness} for \((J_n,\tau)\) consists of four
rational endpoints

\[
W=(a_-,a_+;c_-,c_+)\in\mathbb Q^4.
\tag{3.1}
\]

The first pair is meant to describe an interval around the target that
contains no eigenvalue. The second pair describes an interval that does
contain an eigenvalue. The checker accepts \(W\) precisely when the
following rational conditions hold:

\[
\begin{aligned}
&n\ge1,\quad b_i\ne0\quad(0\le i<n-1),\\
&a_-<\tau<a_+,\qquad N_n(a_-)=N_n(a_+),\\
&c_-<c_+,\qquad N_n(c_-)>N_n(c_+).
\end{aligned}
\tag{3.2}
\]

We write \(\operatorname{Check}(J_n,\tau,W)=\mathrm{true}\) for this
condition. Every comparison in (3.2) is decidable by exact rational
arithmetic and the finite recurrence (2.4)--(2.5). The matrix
admissibility condition is checked, not assumed on the honor of the
witness producer.

Associated to \(W\) are two rational values,

\[
L(W,\tau)=\min\{\tau-a_-,a_+-\tau\},
\qquad
U(W,\tau)=\max\{|c_--\tau|,|c_+-\tau|\}.
\tag{3.3}
\]

The purpose of using two intervals rather than a putative nearest
eigenvalue is that the checker need establish only two finite counting
facts: \emph{absence} in one interval and \emph{presence} in another.
The latter may contain more than one eigenvalue, and need not lie on the
same side of \(\tau\) as the former. Neither fact depends on the method
by which the endpoints were proposed.

\textbf{Theorem 3.1 (soundness of a spectral-distance witness).}
\emph{Let \(n\ge1\), let \(J_n\) be unreduced with rational
coefficients, and let \(\tau\in\mathbb Q\). If the checker accepts a
witness \(W\), then}

\[
0<L(W,\tau)
\le\Delta_n(\tau)
\le U(W,\tau).
\tag{3.4}
\]

\emph{In particular \(\tau\notin\sigma(J_n)\), and the shifted matrix
\(J_n-\tau I\) is invertible.}

\emph{Proof.} The equality \(N_n(a_-)=N_n(a_+)\), together with (2.9),
says that no eigenvalue belongs to \((a_-,a_+]\). Hence every eigenvalue
\(\lambda_j\) lies either at or to the left of \(a_-\), or strictly to
the right of \(a_+\). In the first case,

\[
|\lambda_j-\tau|=\tau-\lambda_j\ge\tau-a_-\ge L(W,\tau),
\]

and in the second,

\[
|\lambda_j-\tau|=\lambda_j-\tau>a_+-\tau\ge L(W,\tau).
\]

Thus the lower bound holds for \emph{every} eigenvalue. Since
\(a_-<\tau<a_+\), both terms in the minimum defining \(L\) are positive.

The strict count drop \(N_n(c_-)>N_n(c_+)\) gives, by (2.9), some
eigenvalue \(\lambda_{j_0}\in(c_-,c_+]\). For any point of that
interval, its distance from \(\tau\) is no greater than the larger
distance of the two endpoints from \(\tau\); consequently,

\[
|\lambda_{j_0}-\tau|
\le\max\{|c_--\tau|,|c_+-\tau|\}
=U(W,\tau).
\]

Taking the minimum over the finite spectrum proves (3.4). Its strict
positive lower bound excludes \(\tau\) from the spectrum and proves
invertibility. \(\square\)

A further inequality, \(L(W,\tau)\le U(W,\tau)\), follows automatically
from (3.4); it need not be presumed of an arbitrary candidate witness. A
defensive verifier may nonetheless check this elementary rational
relation explicitly, both to specify its accepted format and to guard
against mistakes outside the formal core.

\textbf{Corollary 3.2 (resolvent and residual bounds).} \emph{Under the
hypotheses of Theorem 3.1, the Euclidean operator norm satisfies}

\[
\frac1{U(W,\tau)}
\le\|(J_n-\tau I)^{-1}\|_2
\le\frac1{L(W,\tau)}.
\tag{3.5}
\]

\emph{If \(x\) solves \((J_n-\tau I)x=f\) and an approximate solution
\(\widehat x\) has a separately established residual bound
\(\|f-(J_n-\tau I)\widehat x\|_2\le\delta\), then}

\[
\|x-\widehat x\|_2\le\frac{\delta}{L(W,\tau)}.
\tag{3.6}
\]

\emph{Proof.} Since \(J_n\) is real symmetric and \(\tau\) is outside
its spectrum, the spectral theorem yields
\(\|(J_n-\tau I)^{-1}\|_2=1/\Delta_n(\tau)\). Apply (3.4), taking
reciprocals of strictly positive numbers, to obtain (3.5). The error
satisfies \(x-\widehat x=(J_n-\tau I)^{-1}(f-(J_n-\tau I)\widehat x)\);
(3.6) follows by submultiplicativity. \(\square\)

The residual in (3.6) is assumed to have been bounded independently. It
is not an additional object constructed by the spectral-distance
certificate.

\subsection{3.2. Accuracy belongs to the checker
contract}\label{accuracy-belongs-to-the-checker-contract}

Theorem 3.1 establishes two-sided enclosure, but acceptance as defined
in (3.2) does not by itself say that the enclosure meets a caller's
requested tolerance. For a rational \(\varepsilon>0\), define an
\emph{accuracy-qualified} witness by the stronger condition

\[
\operatorname{Check}_{\varepsilon}(J_n,\tau,W)
\quad:\Longleftrightarrow\quad
\operatorname{Check}(J_n,\tau,W)
\ \wedge\ 0<L(W,\tau)\le U(W,\tau)
\ \wedge\ U(W,\tau)-L(W,\tau)\le\varepsilon.
\tag{3.7}
\]

The extra tests involve only rational quantities already present in the
witness. They therefore belong naturally to the \emph{independent
verifier}: a producer's claim that it refined far enough is not evidence
of accuracy until the returned endpoints satisfy (3.7). This separation
is conceptual as well as computational. The producer may be changed or
optimized while the meaning of an accepted certificate remains fixed.

\textbf{Corollary 3.3 (requested-accuracy certification).} \emph{If
\(\operatorname{Check}_{\varepsilon}(J_n,\tau,W)\) holds, then}

\[
0<L(W,\tau)\le\Delta_n(\tau)\le U(W,\tau),
\qquad U(W,\tau)-L(W,\tau)\le\varepsilon.
\tag{3.8}
\]

\emph{Proof.} The distance enclosure is Theorem 3.1; the width bound is
an explicitly checked rational inequality. \(\square\)

For an accuracy-requesting computational interface, the appropriate
acceptance criterion is (3.7), rather than (3.2) alone. The independent
rational width comparison is part of the certificate semantics, not an
unverified assertion by its producer. The width comparison is
deliberately made by the verifier from the returned rational endpoints,
rather than accepted as a producer assertion.

\subsection{3.3. Bracketing the distance
itself}\label{bracketing-the-distance-itself}

The preceding witness may have quite general intervals. A particularly
economical producer uses \emph{symmetric intervals about the target}.
For \(r>0\), put

\[
D_\tau(r):=N_n(\tau-r)-N_n(\tau+r)
=\#\{j:\tau-r<\lambda_j\le\tau+r\}.
\tag{3.9}
\]

By Theorem 2.1, \(D_\tau(r)\ge0\), and it is nondecreasing in \(r\). Its
exact value can be found by two rational Sturm-count evaluations. In
particular,

\[
\begin{array}{rcl}
D_\tau(r)=0&\Longrightarrow&\Delta_n(\tau)\ge r,\\
D_\tau(r)>0&\Longrightarrow&\Delta_n(\tau)\le r.
\end{array}
\tag{3.10}
\]

The weak inequalities in (3.10) are intentional. An eigenvalue at the
left endpoint \(\tau-r\) belongs to neither the half-open interval nor
its count, but its distance from \(\tau\) is exactly \(r\); no erroneous
strict inequality is inferred.

Suppose we have rational radii \(0<\ell<u\) satisfying

\[
D_\tau(\ell)=0,
\qquad D_\tau(u)>0.
\tag{3.11}
\]

They determine a witness of the type (3.1) by setting

\[
(a_-,a_+)=(\tau-\ell,\tau+\ell),
\qquad
(c_-,c_+)=(\tau-u,\tau+u).
\tag{3.12}
\]

The corresponding endpoints simplify to

\[
L(W,\tau)=\ell,
\qquad U(W,\tau)=u.
\tag{3.13}
\]

This is why we bisect radii rather than eigenvalues: the \emph{width of
the distance certificate} is exactly \(u-\ell\).

\textbf{Proposition 3.4 (count-preserving radial bisection).}
\emph{Assume (3.11) and put \(m=(\ell+u)/2\). Define}

\[
(\ell',u')=
\begin{cases}
(\ell,m),&D_\tau(m)>0,\\
(m,u),&D_\tau(m)=0.
\end{cases}
\tag{3.14}
\]

\emph{Then \(0<\ell'<u'\), \(D_\tau(\ell')=0\), \(D_\tau(u')>0\), and}

\[
u'-\ell'=\frac{u-\ell}{2}.
\tag{3.15}
\]

\emph{After exactly \(k\) successive steps,}

\[
u_k-\ell_k=\frac{u_0-\ell_0}{2^k}.
\tag{3.16}
\]

\emph{Proof.} Since \(0<\ell<m<u\), either branch of (3.14) preserves
strict ordering and positivity. If \(D_\tau(m)>0\), the new outer
endpoint is occupied and the old inner endpoint remains empty. Otherwise
the count difference is zero (it cannot be negative, by (2.9)), so the
new inner endpoint is empty and the outer endpoint remains occupied.
Both branches have width \((u-\ell)/2\); induction gives (3.16).
\(\square\)

No special branch is required when \(m\) itself belongs to the spectrum.
The strict spectral count is defined at such a point, and the half-open
convention determines exactly which radius statement remains valid.

It follows from (3.16) that any pre-existing valid pair can be refined
to width at most a given \(\varepsilon>0\). But (3.16) is an exact width
identity, not a bit-complexity estimate. A terminating implementation
may still take many unnecessary steps, and the intermediate rational
numerators and denominators may grow. Numerical usefulness therefore
requires both a sensible finite stopping budget and measurements of
exact-arithmetic cost.

\subsection{3.4. The three-site certificate
revisited}\label{the-three-site-certificate-revisited}

For the matrix in the introduction, \(q_0(x)=1\), \(q_1(x)=x\),
\(q_2(x)=x^2-1\), and \(q_3(x)=x^3-2x\). Direct rational evaluation
gives:

\begin{longtable}[]{@{}lcr@{}}
\toprule\noalign{}
Test energy \(x\) & Signs of \((q_0,q_1,q_2,q_3)(x)\) & \(N_3(x)\) \\
\midrule\noalign{}
\endhead
\bottomrule\noalign{}
\endlastfoot
\(139/100\) & \((+,+,+,-)\) & 1 \\
\(141/100\) & \((+,+,+,-)\) & 1 \\
\(71/50\) & \((+,+,+,+)\) & 0 \\
\end{longtable}

Hence the exclusion condition holds on \((139/100,141/100]\), and the
occupied condition holds on \((141/100,71/50]\). With \(\tau=7/5\),
equations (3.3) and (3.4) give

\[
L=\frac1{100},\qquad U=\frac1{50},\qquad
\frac1{100}\le\Delta_3(7/5)\le\frac1{50}.
\tag{3.17}
\]

The certified interval has width \(1/100\). The consequent resolvent
enclosure is

\[
50\le\|(J-\tfrac75 I)^{-1}\|_2\le100.
\tag{3.18}
\]

This example demonstrates the meaning of an accepted certificate; it is
not a benchmark for the asymptotic efficiency of a producer. The
relevant claim is a theorem about actual eigenvalues, not merely the
successful execution of an arithmetic test.

\subsection{3.5. The separation of production and
verification}\label{the-separation-of-production-and-verification}

The mathematical soundness argument above does not use an eigenvalue
approximation, an iteration count, or the internal structure of the
algorithm that proposed \(W\). A producer is free to exploit exact
bisection, continuant perturbation bounds, or other methods; acceptance
depends only on the data and the finite checks in (3.2) and, for
requested accuracy, (3.7). This is the intended trust division of the
accompanying formal development.

This division of trust does not by itself verify serialization, the
surrounding runtime, or the cost of exact arithmetic. Its significance
is more precise: \textbf{the implication from accepted rational evidence
to a true spectral-distance bound can be proved independently of its
producer.}

The next section will explain how to construct the initial empty and
occupied radii without knowing an eigenvalue in advance. The present
sections have established what evidence those radii must provide---and
what, exactly, may be concluded when it is accepted.

\section{4. Constructing the
certificate}\label{constructing-the-certificate}

The preceding section asks for two rational radii: a positive radius
whose neighborhood is empty of eigenvalues, and a larger radius whose
neighborhood is occupied. Neither radius may be supplied by an
approximate eigenvalue whose accuracy is taken on trust. In this section
we construct both directly from the Jacobi coefficients, then show how
finitely many rational Sturm tests yield an arbitrarily narrow distance
enclosure.

Throughout, the data satisfy (2.1), the target \(\tau\) is rational, and
\(\varepsilon>0\) is rational. We begin with the exact test
\(q_n(\tau)=0\). By (2.6), a zero is already a certificate that \(\tau\)
is an eigenvalue. The nonzero case is more interesting: continuity
guarantees some neighborhood free of eigenvalues, but a terminating
exact algorithm needs to \emph{calculate} a positive radius for that
neighborhood.

\subsection{4.1. A computable root-free
neighborhood}\label{a-computable-root-free-neighborhood}

Set \(C_0=1\), \(C_1=|\tau-d_0|+1\) and \(S_0=0\), \(S_1=1\). For
\(0\le k\le n-2\) define the rational sequences

\[
\begin{aligned}
C_{k+2}&=(|\tau-d_{k+1}|+1)C_{k+1}+b_k^2C_k,\\
S_{k+2}&=C_{k+1}+|\tau-d_{k+1}|S_{k+1}+b_k^2S_k.
\end{aligned}
\tag{4.1}
\]

Both sequences are nonnegative and are computed by the same forward
traversal of the coefficients as the characteristic continuant. The two
sequences have different roles: \(C_k\) bounds the value of \(q_k\)
throughout a unit neighborhood of \(\tau\), while \(S_k\) bounds its
variation with the energy. These are intentionally coarse majorants, not
estimates of the nearest eigenvalue.

\textbf{Lemma 4.1 (continuant perturbation bound).} \emph{For every
\(0\le k\le n\) and every real \(x\) satisfying \(|x|\le1\),}

\[
|q_k(\tau+x)|\le C_k,
\qquad
|q_k(\tau+x)-q_k(\tau)|\le |x|S_k.
\tag{4.2}
\]

\emph{Proof.} We prove both inequalities simultaneously by induction
through the two preceding continuants. For \(k=0\), \(q_0=1\) and both
assertions are immediate. For \(k=1\),

\[
|\tau+x-d_0|\le|\tau-d_0|+|x|\le C_1,
\qquad
|q_1(\tau+x)-q_1(\tau)|=|x|=|x|S_1.
\]

Suppose (4.2) holds at indices \(k\) and \(k+1\). Write
\(a=\tau-d_{k+1}\) and \(B=b_k^2\ge0\). The recurrence (2.5) and the
triangle inequality give

\[
\begin{aligned}
|q_{k+2}(\tau+x)|
&\le (|a|+|x|)|q_{k+1}(\tau+x)|
       +B|q_k(\tau+x)|\\
&\le (|a|+1)C_{k+1}+BC_k=C_{k+2}.
\end{aligned}
\]

For the difference, subtract the recurrence at \(\tau\) from that at
\(\tau+x\):

\[
\begin{aligned}
q_{k+2}(\tau+x)-q_{k+2}(\tau)
={}&xq_{k+1}(\tau+x)\\
&+a\bigl(q_{k+1}(\tau+x)-q_{k+1}(\tau)\bigr)\\
&-B\bigl(q_k(\tau+x)-q_k(\tau)\bigr).
\end{aligned}
\]

Applying the induction hypotheses to these three terms yields

\[
|q_{k+2}(\tau+x)-q_{k+2}(\tau)|
\le |x|\bigl(C_{k+1}+|a|S_{k+1}+BS_k\bigr)
=|x|S_{k+2}.
\]

This completes the induction. The argument is over \emph{real} shifts
\(x\), even though its coefficients and majorants are rational.
\(\square\)

Assume now that \(q_n(\tau)\ne0\), and define

\[
\boxed{
 r_{\mathrm{in}}=
 \min\left\{\frac12,
 \frac{|q_n(\tau)|}{2(S_n+1)}\right\}.}
\tag{4.3}
\]

Every quantity in (4.3) is rational and computed from the supplied
matrix and target. In particular, \(S_n\ge0\) and \(q_n(\tau)\ne0\)
imply \(r_{\mathrm{in}}>0\).

\textbf{Proposition 4.2 (constructive spectral exclusion).} \emph{If
\(q_n(\tau)\ne0\), then \(q_n(t)\ne0\) for every real \(t\) with
\(|t-\tau|\le r_{\mathrm{in}}\). Consequently,}

\[
N_n(\tau-r_{\mathrm{in}})
=N_n(\tau+r_{\mathrm{in}}),
\qquad
\Delta_n(\tau)\ge r_{\mathrm{in}}>0.
\tag{4.4}
\]

\emph{Proof.} Write \(t=\tau+x\) with \(|x|\le r_{\mathrm{in}}\le1/2\).
From Lemma 4.1 and the definition of the radius,

\[
\begin{aligned}
|q_n(t)-q_n(\tau)|
&\le |x|S_n
\le r_{\mathrm{in}}S_n\\
&\le \frac{|q_n(\tau)|S_n}{2(S_n+1)}
\le\frac{|q_n(\tau)|}{2}.
\end{aligned}
\]

The reverse triangle inequality therefore gives

\[
|q_n(t)|\ge |q_n(\tau)|-|q_n(t)-q_n(\tau)|
\ge\frac{|q_n(\tau)|}{2}>0.
\tag{4.5}
\]

By (2.6), no eigenvalue lies in the closed real interval
\([\tau-r_{\mathrm{in}},\tau+r_{\mathrm{in}}]\). Its half-open
subinterval contains no eigenvalue either, and (2.9) proves the asserted
equality of Sturm counts. The distance bound follows. \(\square\)

The closed-interval conclusion is a useful strength of this
construction. It prevents any hidden appeal to a rational-only test: the
exact rational radius excludes \emph{all real spectral roots} in its
neighborhood.

\subsection{4.2. An occupied outer neighborhood without an eigenvalue
search}\label{an-occupied-outer-neighborhood-without-an-eigenvalue-search}

To bound the distance from above, it is enough to find a rational
\(R>r_{\mathrm{in}}\) such that \((\tau-R,\tau+R]\) contains an
eigenvalue. We need not approximate, select, or isolate that eigenvalue.

For the next formula only, adopt the harmless finite-array convention
\(b_{n-1}=0\). This is \emph{padding}, not an additional off-diagonal of
\(J_n\). Define the rational budget

\[
H_n=2\sum_{k=0}^{n-1}\bigl(|d_k|+b_k^2+2\bigr),
\qquad
\boxed{R=r_{\mathrm{in}}+|\tau|+H_n+2.}
\tag{4.6}
\]

The budget is deliberately larger than needed; it has the advantage of
an elementary certificate directly in terms of the continuants. In
particular, \(R>r_{\mathrm{in}}\).

\textbf{Lemma 4.3 (explicit occupied outer radius).} \emph{For
\(n\ge1\), the radius (4.6) satisfies}

\[
N_n(\tau+R)=0,
\qquad N_n(\tau-R)>0.
\tag{4.7}
\]

\emph{Proof.} Set \(x_+=\tau+R\). The coefficient budget guarantees

\[
x_+-d_0>1,
\qquad
x_+-d_{k+1}>1+b_k^2
\quad(0\le k\le n-2).
\tag{4.8}
\]

We claim that the continuants evaluated at \(x_+\) are positive and
strictly increasing from \(q_0\) onward. Indeed, \(q_0(x_+)=1\) and
\(q_1(x_+)=x_+-d_0>1\). If \(0<q_k(x_+)<q_{k+1}(x_+)\), then the
recurrence and (4.8) imply

\[
\begin{aligned}
q_{k+2}(x_+)-q_{k+1}(x_+)
&=(x_+-d_{k+1}-1)q_{k+1}(x_+)-b_k^2q_k(x_+)\\
&>b_k^2\bigl(q_{k+1}(x_+)-q_k(x_+)\bigr)\ge0.
\end{aligned}
\]

Thus the entire sign list is positive, and \(N_n(x_+)=0\).

On the other side put \(x_-=\tau-R\). The same budget ensures
\(x_-<d_0\). Thus the first two continuants have opposite signs:

\[
q_0(x_-)=1>0,
\qquad q_1(x_-)=x_--d_0<0.
\]

Whatever signs or zeros occur later in the list, they cannot erase this
first variation. Hence \(N_n(x_-)\ge1\). \(\square\)

The asymmetry of the proof is instructive. On the right we establish
positivity of every continuant. On the left we need only the first sign
change, not a second positivity or alternation induction. Combined with
(2.9), (4.7) certifies that the outer neighborhood is occupied. The
large budget buys a simple proof of existence, not a claim of numerical
sharpness.

\subsection{4.3. Exact refinement and a finite stopping
rule}\label{exact-refinement-and-a-finite-stopping-rule}

Propositions 4.2 and Lemma 4.3 initialize the radial pair

\[
\ell_0=r_{\mathrm{in}},\qquad u_0=R,
\qquad 0<\ell_0<u_0,
\tag{4.9}
\]

with \(D_\tau(\ell_0)=0\) and \(D_\tau(u_0)>0\), where \(D_\tau\) is the
count difference (3.9). The bisection step (3.14) therefore applies as
often as desired. Each step uses only rational arithmetic and two exact
Sturm-count evaluations; its validity does not depend on whether a
midpoint happens to be an eigenvalue.

To stop at tolerance \(\varepsilon>0\), it suffices to choose an integer
\(k\ge0\) with

\[
2^k\ge\frac{u_0-\ell_0}{\varepsilon}.
\tag{4.10}
\]

The exact width identity (3.16) then gives

\[
0<u_k-\ell_k
=\frac{u_0-\ell_0}{2^k}
\le\varepsilon.
\tag{4.11}
\]

The existence of such a finite integer is elementary and constructive.
For example, the deliberately conservative choice

\[
k=\left\lceil\frac{u_0-\ell_0}{\varepsilon}\right\rceil+1
\tag{4.12}
\]

suffices by \(k\le2^k\). This proves finite termination without any
floating-point logarithm. It is, however, a poor work budget when
\(\varepsilon\) is small: the \emph{mathematically natural} choice is
the least \(k\) satisfying (4.10), whose size grows logarithmically with
\((u_0-\ell_0)/\varepsilon\). The refinement theorem is independent of
which certified sufficient budget is used. An exact implementation must
not confuse the convenience of (4.12) in a totality proof with practical
iteration complexity.

Convert the final radii to the witness

\[
W_k=(\tau-\ell_k,\tau+\ell_k;
      \tau-u_k,\tau+u_k).
\tag{4.13}
\]

By (3.11)--(3.13) and (4.11), this witness passes the exact spectral
checker and has \(L(W_k,\tau)=\ell_k\), \(U(W_k,\tau)=u_k\), and
\(U-L\le\varepsilon\). All of these properties are rational relations
available for independent checking.

\textbf{Theorem 4.4 (constructive rational spectral-distance
certification).} \emph{Fix an unreduced rational \(n\)-site Jacobi
matrix with \(n\ge1\), a rational target \(\tau\), and a rational
tolerance \(\varepsilon>0\). A finite exact procedure distinguishes the
following cases:}

\begin{enumerate}
\def\labelenumi{\arabic{enumi}.}
\tightlist
\item
  \emph{If \(q_n(\tau)=0\), it certifies that \(\det(\tau I_n-J_n)=0\),
  hence that the shifted matrix is singular.}
\item
  \emph{If \(q_n(\tau)\ne0\), it constructs a rational witness \(W\) for
  which}
\end{enumerate}

\[
\operatorname{Check}_{\varepsilon}(J_n,\tau,W)
\quad\text{and}\quad
0<L(W,\tau)\le\Delta_n(\tau)\le U(W,\tau),
\quad U(W,\tau)-L(W,\tau)\le\varepsilon.
\tag{4.14}
\]

\emph{The soundness of the second conclusion depends only on
independently checkable properties of \(W\), not on trusting the
construction that produced it.}

\emph{Proof.} Compute \(q_n(\tau)\) exactly by (2.4)--(2.5). In the zero
case, (2.6) establishes the stated determinant identity. Otherwise,
compute (4.1), (4.3), and (4.6), obtaining an empty positive inner
radius and an occupied larger outer radius by Proposition 4.2 and Lemma
4.3. Choose the finite integer \(k\) as in (4.10), perform \(k\) exact
count-preserving bisections, and form (4.13). Proposition 3.4 proves
preservation of the empty/occupied invariants and (4.11) proves the
requested width. Equations (3.2), (3.7), and (3.13) establish acceptance
by the accuracy-qualified checker. Theorem 3.1 supplies the spectral
conclusion independently of the producer. \(\square\)

The theorem is a statement about \emph{exact rational input}. It does
not presume rational eigenvalues. Nor does it claim a verified
floating-point front end or an efficient bit-complexity bound. The role
of the producer is to propose suitable rational evidence; the checker
need only recompute the finite predicates that give that evidence its
meaning.

\subsection{4.4. What the initial radius looks like in the three-site
example}\label{what-the-initial-radius-looks-like-in-the-three-site-example}

For the matrix of Section 1 and \(\tau=7/5\), the continuant at the
target is

\[
q_3(7/5)=(7/5)^3-2(7/5)=-\frac7{125}.
\]

The rational majorants from (4.1) are

\[
(C_0,C_1,C_2,C_3)
=\left(1,\frac{12}{5},\frac{169}{25},\frac{2328}{125}\right),
\qquad
(S_0,S_1,S_2,S_3)
=\left(0,1,\frac{19}{5},\frac{327}{25}\right).
\tag{4.15}
\]

It follows that

\[
r_{\mathrm{in}}
=\min\left\{\frac12,
 \frac{7/125}{2(327/25+1)}\right\}
=\frac7{3520}>0.
\tag{4.16}
\]

This is substantially smaller than the actual separation near
\(\sqrt2\), and intentionally so: it is obtained without knowing
\(\sqrt2\) or invoking a numerical eigenvalue approximation. With the
padded convention \(b_2=0\), the outer budget is \(H_3=16\), giving
\(R=7/3520+97/5\). Exact bisection then narrows the rather wide initial
enclosure to the requested rational width. The much sharper
hand-selected witness in Section 3 illustrates what the independent
checker can validate; it is not claimed to coincide with these
particular producer iterates.

The constructive theorem proves that suitable evidence can always be
found in the nonsingular rational case. Its deliberately coarse
initialization and stopping budget leave a separate question about
efficient execution, addressed by the subsequent computational analysis
rather than by altering the truth of the certificate.

\section{5. Conditioning under represented
uncertainty}\label{conditioning-under-represented-uncertainty}

A certificate concerns an exact pair \((J_n,\tau)\). Its practical value
extends beyond that pair if the matrix entries or target energy are
known only within independently certified bounds. The distinction is
important: the rational producer above does not itself validate an
external approximation to an irrational target, and it does not certify
uncertain matrix coefficients. The next result says precisely what
additional bounds suffice to \emph{transport} the already certified
conclusion.

\subsection{5.1. Spectral distance is
Lipschitz}\label{spectral-distance-is-lipschitz}

Write \(\Delta(A,t)=\min_{\lambda\in\sigma(A)}|\lambda-t|\) for a real
symmetric \(n\times n\) matrix \(A\), where \(n\ge1\). All operator
norms below are induced by the Euclidean vector norm.

\textbf{Theorem 5.1 (spectral-distance transport).} \emph{If \(A\) and
\(B\) are real symmetric matrices of the same positive dimension and
\(s,t\in\mathbb R\), then}

\[
\boxed{
|\Delta(A,t)-\Delta(B,s)|
\le\|A-B\|_2+|t-s|.}
\tag{5.1}
\]

\emph{Proof.} Arrange the eigenvalues of \(A\) and \(B\) in
nondecreasing order, \(\lambda_1(A)\le\cdots\le\lambda_n(A)\) and
likewise for \(B\). Weyl's eigenvalue perturbation inequality gives

\[
|\lambda_i(A)-\lambda_i(B)|\le\|A-B\|_2
\qquad (1\le i\le n).
\]

Therefore for each index,

\[
\bigl||\lambda_i(A)-t|-|\lambda_i(B)-s|\bigr|
\le\|A-B\|_2+|t-s|.
\]

Taking a minimizing index first for \(A\) and then for \(B\) yields the
two opposite one-sided estimates; together they give (5.1). \(\square\)

For a nominal rational matrix \(J_n\) and rational target \(\tau\),
suppose an accepted witness yields \(L\le\Delta(J_n,\tau)\le U\). Let a
real symmetric matrix \(\widetilde J_n\) and real target
\(\widetilde\tau\) satisfy \emph{independently justified} bounds

\[
\|\widetilde J_n-J_n\|_2\le\eta,
\qquad
|\widetilde\tau-\tau|\le\rho,
\qquad \eta,\rho\ge0.
\tag{5.2}
\]

Applying Theorem 5.1 gives

\[
\boxed{
\max\{0,L-\eta-\rho\}
\le\Delta(\widetilde J_n,\widetilde\tau)
\le U+\eta+\rho.}
\tag{5.3}
\]

The loss of lower-bound positivity is meaningful. If \(L\le\eta+\rho\),
the original certificate no longer excludes singularity for every pair
in the uncertainty set. If instead \(L>\eta+\rho\), every such shifted
system is invertible, and

\[
\frac1{U+\eta+\rho}
\le
\|(\widetilde J_n-\widetilde\tau I)^{-1}\|_2
\le
\frac1{L-\eta-\rho}.
\tag{5.4}
\]

This is a robustness conclusion, not a new spectral-perturbation
inequality. The contribution of the certificate is that the initial
margin \(L\) was established by exact finite evidence.

\subsection{5.2. Converting coefficient enclosures to an operator-norm
bound}\label{converting-coefficient-enclosures-to-an-operator-norm-bound}

For tridiagonal uncertainty, the norm input in (5.2) can itself be
bounded by elementary data. Suppose \(\widetilde J_n\) has real diagonal
coefficients \(\widetilde d_i\) and real off-diagonal coefficients
\(\widetilde b_i\), and nonnegative rational numbers
\(\alpha_i,\beta_i\) have been certified to satisfy

\[
|\widetilde d_i-d_i|\le\alpha_i
\quad(0\le i<n),
\qquad
|\widetilde b_i-b_i|\le\beta_i
\quad(0\le i<n-1).
\tag{5.5}
\]

With the endpoint convention \(\beta_{-1}=\beta_{n-1}=0\), set

\[
\eta_{\mathrm{row}}
:=\max_{0\le i<n}\bigl(\alpha_i+\beta_{i-1}+\beta_i\bigr).
\tag{5.6}
\]

\textbf{Corollary 5.2 (entrywise uncertainty budget).} \emph{Under
(5.5),}

\[
\|\widetilde J_n-J_n\|_2\le\eta_{\mathrm{row}},
\tag{5.7}
\]

\emph{so (5.3)--(5.4) hold with \(\eta=\eta_{\mathrm{row}}\).}

\emph{Proof.} The symmetric difference matrix has at most three nonzero
entries per row. By (5.5), the sum of the absolute values in row \(i\)
is bounded by \(\alpha_i+\beta_{i-1}+\beta_i\). Its induced infinity
norm is therefore at most \(\eta_{\mathrm{row}}\). Symmetry implies
equality of the induced one- and infinity-norm bounds, and the standard
inequality \(\|E\|_2\le\sqrt{\|E\|_1\|E\|_\infty}\) yields (5.7).
\(\square\)

Notice that \(\widetilde J_n\) need not be unreduced. The nonzero
off-diagonal hypothesis is needed for the \emph{nominal exact Sturm
certificate}; the transported bound applies to any real symmetric
perturbation satisfying its norm estimate.

\subsection{5.3. The safe-solve
consequence}\label{the-safe-solve-consequence}

Let \(\widetilde A=\widetilde J_n-\widetilde\tau I\), and suppose
\(L>\eta+\rho\). If \(\widehat x\) is a proposed solution of
\(\widetilde A x=f\) and a separate computation certifies a residual
bound

\[
\|f-\widetilde A\widehat x\|_2\le\delta,
\]

then the exact solution exists uniquely and satisfies

\[
\boxed{
\|x-\widehat x\|_2
\le\frac{\delta}{L-\eta-\rho}.}
\tag{5.8}
\]

Indeed, apply the upper inverse-norm bound in (5.4) to the residual
equation. In this sense a spectral certificate answers three connected
questions: whether the shifted system can be solved uniquely, how much
the inverse may amplify a residual, and how much perturbation the
certified separation can withstand.

The external enclosures in (5.2) and (5.5) must themselves be valid.
Nothing in the preceding argument promotes an unverified decimal
approximation into a certified real interval. Likewise, the certificate
does not generate the residual \(\delta\). Its role is to supply a
rigorous spectral margin that can be composed with those independently
established facts.

\section{6. Arithmetic resonance and the resolution of spectral
separation}\label{arithmetic-resonance-and-the-resolution-of-spectral-separation}

The certificate of Section 3 does not require any arithmetic structure
in the spectrum. For a constant Jacobi chain, however, its distance to a
prescribed energy can be expressed in terms of rational approximation to
a phase. This makes a useful stress test: the matrix entries are
elementary, yet the separation to be certified can be exceptionally
small. Our purpose is to identify that scale and to distinguish a
requirement of one explicitly specified information model from the
computational cost of the exact Sturm certificate itself.

\subsection{6.1. The Dirichlet spectral
mesh}\label{the-dirichlet-spectral-mesh}

For an integer \(q\ge2\), let \(T_{q-1}(c,\beta)\) be the
\((q-1)\times(q-1)\) Jacobi matrix with constant real diagonal \(c\) and
nonzero constant real off-diagonal \(\beta\). Its Dirichlet eigenvalues
are

\[
\lambda_{q,k}=c+2\beta\cos\frac{\pi k}{q},
\qquad 1\le k\le q-1.
\tag{6.1}
\]

Indeed, the vector with components \(v_j=\sin(\pi kj/q)\),
\(1\le j\le q-1\), obeys the boundary conditions \(v_0=v_q=0\), and the
identity \(v_{j-1}+v_{j+1}=2\cos(\pi k/q)v_j\) gives the eigenvalue
equation. The values in (6.1) are distinct because cosine is strictly
decreasing on \((0,\pi)\) and \(\beta\ne0\).

Fix throughout this section an \emph{irrational} \(\alpha\in(0,1)\) and
the real target

\[
\tau_\alpha=c+2\beta\cos(\pi\alpha).
\tag{6.2}
\]

Irrationality ensures that \(\tau_\alpha\) is not an eigenvalue of any
of these finite chains. With
\(\Delta_q:=\Delta(T_{q-1}(c,\beta),\tau_\alpha)\), equation (6.1)
yields the exact finite identity

\[
\Delta_q=2|\beta|\min_{1\le k\le q-1}
  \left|\cos(\pi\alpha)-\cos\frac{\pi k}{q}\right|>0.
\tag{6.3}
\]

No limit has been taken in (6.3). Since the target is nonspectral, the
ordinary Euclidean operator norm of the inverse satisfies

\[
\left\|\bigl(T_{q-1}(c,\beta)-\tau_\alpha I\bigr)^{-1}\right\|_2
=\Delta_q^{-1}.
\tag{6.4}
\]

The following comparison identifies the arithmetic feature behind this
inverse norm. Write

\[
\|x\|_{\mathbb Z}:=\min_{m\in\mathbb Z}|x-m|,
\qquad \delta_q:=\|q\alpha\|_{\mathbb Z},
\qquad s_\alpha:=\min\left\{\sin\frac{\pi\alpha}{2},
\cos\frac{\pi\alpha}{2}\right\}>0.
\tag{6.5}
\]

\textbf{Proposition 6.1 (explicit phase-to-distance comparison).}
\emph{For fixed \(\alpha\in(0,1)\), irrational, and \(\beta\ne0\), every
integer}

\[
q>\frac{1}{\min\{\alpha,1-\alpha\}}
\tag{6.6}
\]

\emph{satisfies}

\[
\boxed{
\frac{4|\beta|s_\alpha}{q}\,\delta_q
\ \le\ \Delta_q\ \le\
\frac{2\pi|\beta|}{q}\,\delta_q.}
\tag{6.7}
\]

\emph{Proof.} For \(t\in[0,1]\), the cosine difference identity gives

\[
|\cos(\pi t)-\cos(\pi\alpha)|
=2\sin\!\left(\frac{\pi(t+\alpha)}2\right)
 \left|\sin\!\left(\frac{\pi(t-\alpha)}2\right)\right|.
\]

The first sine is positive and, by concavity of sine on \([0,\pi]\), is
at least the smaller of its endpoint values as \(t\) ranges over
\([0,1]\). Those values are \(\sin(\pi\alpha/2)\) and
\(\cos(\pi\alpha/2)\), so the factor is at least \(s_\alpha\). The
elementary estimate \(|\sin u|\ge2|u|/\pi\) for \(|u|\le\pi/2\) then
gives

\[
|\cos(\pi t)-\cos(\pi\alpha)|\ge2s_\alpha|t-\alpha|.
\tag{6.8}
\]

The mean-value theorem supplies the converse inequality
\(|\cos(\pi t)-\cos(\pi\alpha)|\le\pi|t-\alpha|\).

Condition (6.6) puts \(q\alpha\) strictly between \(1\) and \(q-1\);
hence a nearest integer to \(q\alpha\) is among \(1,\ldots,q-1\).
Therefore \(\min_{1\le k\le q-1}|k/q-\alpha|=\delta_q/q\). Apply the two
cosine inequalities to each \(k/q\), take the minimum, and multiply by
\(2|\beta|\). \(\square\)

The constants in (6.7) are intentionally explicit rather than sharp. In
particular, they depend on the fixed phase \(\alpha\) and coupling
\(\beta\), but not on \(q\). Combining (6.4) and (6.7) gives

\[
\boxed{
\frac{q}{2\pi|\beta|\delta_q}
\ \le\
\left\|\bigl(T_{q-1}-\tau_\alpha I\bigr)^{-1}\right\|_2
\ \le\
\frac{q}{4|\beta|s_\alpha\delta_q}.}
\tag{6.9}
\]

Here and below, \(T_{q-1}\) abbreviates \(T_{q-1}(c,\beta)\). The
familiar shorthand \(\Delta_q\asymp_{\alpha,\beta}\delta_q/q\) refers
precisely to a comparison such as (6.7), \emph{with the phase held
fixed}. It is not a uniform estimate as \(\alpha\) approaches an
endpoint of \((0,1)\).

\subsection{6.2. A sharp threshold in a specified interval-information
model}\label{a-sharp-threshold-in-a-specified-interval-information-model}

There are many ways to certify spectral separation. To make a precise
statement about required resolution, we first fix one particular
reporting format. Let a real symmetric \(n\times n\) matrix have
eigenvalues \(\lambda_1,\ldots,\lambda_n\), let \(\tau\notin\sigma(J)\),
and write \(d=\Delta(J,\tau)>0\). Suppose a reporting procedure returns,
for each eigenvalue, a \emph{closed symmetric interval}

\[
I_j=[m_j-\varepsilon,m_j+\varepsilon]
\quad\text{with}\quad
|m_j-\lambda_j|\le\varepsilon.
\tag{6.10}
\]

The centers \(m_j\) are allowed to be any real values satisfying these
certified error bounds. Consider the exclusion rule which declares
nonresonance when \(\tau\notin\bigcup_j I_j\). The rule is sound
whenever it reports exclusion, but the error budget alone need not force
it to succeed for every admissible set of centers.

\textbf{Proposition 6.2 (sharp uniform-exclusion threshold).}
\emph{Under the information model (6.10), the rule excludes \(\tau\) for
\textbf{every} admissible choice of centers if and only if}

\[
\boxed{2\varepsilon<d.}
\tag{6.11}
\]

\emph{Proof.} For any admissible center,
\(|m_j-\tau|\ge|\lambda_j-\tau|-|m_j-\lambda_j|\ge d-\varepsilon\). If
\(2\varepsilon<d\), then \(|m_j-\tau|>\varepsilon\) for each \(j\),
which is exactly \(\tau\notin I_j\).

Conversely, suppose \(2\varepsilon\ge d\). Choose an eigenvalue
\(\lambda_{j_*}\) at distance \(d\) from \(\tau\), and move its reported
center toward \(\tau\) by \(\min\{d,\varepsilon\}\). The center remains
within \(\varepsilon\) of \(\lambda_{j_*}\), while its distance from
\(\tau\) is \(\max\{0,d-\varepsilon\}\le\varepsilon\). Thus
\(\tau\in I_{j_*}\) for this admissible report. At equality
\(2\varepsilon=d\), the target can lie at the endpoint of the closed
interval; endpoint contact is not exclusion. \(\square\)

Take the absolute half-width \(\varepsilon=2^{-p}\), with \(p\) a
nonnegative integer. For sufficiently small fixed \(d>0\), the least
integer \(p\) guaranteeing uniform exclusion by this \emph{particular}
rule is

\[
p_*(d)=\left\lfloor\log_2\frac2d\right\rfloor+1,
\qquad\text{provided }d<2,
\tag{6.12}
\]

because strict exclusion requires \(p>\log_2(2/d)\). If \(d\ge2\), the
smallest admissible \(p\) is instead
\(\max\{0,\lfloor\log_2(2/d)\rfloor+1\}\); thus the formula can be
written uniformly with that maximum. This \(p\) is an \emph{absolute
interval half-width exponent}. It is not a floating-point mantissa
length, a bit-operation count, or a lower bound on algorithms with
different information or certificate formats.

Applying (6.12) to the constant-chain distance in (6.7), and using
\(\Delta_q\to0\), yields the next conclusion.

\textbf{Corollary 6.3 (arithmetic scale of symmetric-interval
resolution).} \emph{For fixed irrational \(\alpha\in(0,1)\), fixed
\(\beta\ne0\), and \(q\to\infty\), let \(p_q\) be the least half-width
exponent for which every admissible report (6.10) excludes
\(\tau_\alpha\) from the spectrum of \(T_{q-1}(c,\beta)\). Then}

\[
\boxed{
 p_q=\log_2\!\left(\frac{q}{\|q\alpha\|_{\mathbb Z}}\right)
       +O_{\alpha,\beta}(1).}
\tag{6.13}
\]

\emph{Proof.} The two inequalities in (6.7) place \(\Delta_q\) between
fixed positive multiples of \(\delta_q/q\). Taking
\(\log_2(2/\Delta_q)\) changes \(\log_2(q/\delta_q)\) by a bounded
quantity depending only on \(\alpha\) and \(\beta\). Equation (6.12)
introduces at most one further unit through integer rounding.
\(\square\)

The qualifier \emph{uniform over admissible reports} cannot be omitted.
The proposition describes how accurately a full list of symmetric
eigenvalue enclosures must be reported to guarantee separation
independently of their centers. Our Sturm certificate does not report
such a list: it tests the occupancy of two intervals centered at the
target. Therefore (6.13) is an illustration of arithmetic precision
demand in one information model, not a statement that every
spectral-distance algorithm must use \(p_q\) bits or perform at least
\(p_q\) operations.

\subsection{6.3. Two Diophantine regimes}\label{two-diophantine-regimes}

If \(\alpha\) is \emph{badly approximable}, then there is a constant
\(\kappa_\alpha>0\) such that

\[
\|q\alpha\|_{\mathbb Z}\ge\frac{\kappa_\alpha}{q}
\qquad(q\ge1).
\tag{6.14}
\]

By (6.7), \(\Delta_q\) is bounded below by a constant multiple of
\(q^{-2}\), and (6.13) gives \(p_q\le2\log_2q+O_{\alpha,\beta}(1)\).
Conversely, the denominators of continued-fraction convergents of any
irrational \(\alpha\) include an unbounded sequence for which
\(\|q\alpha\|_{\mathbb Z}<1/q\). Along this sequence, (6.7) gives
\(\Delta_q=O_\beta(q^{-2})\) and the threshold is at least
\(2\log_2q-O_\beta(1)\). Thus for a badly approximable phase the largest
resolution demand along the denominators has quadratic spectral scale,
up to phase-dependent constants.

Quadratic irrational phases, including \(\alpha=\sqrt2-1\), are familiar
examples of this regime. For this particular phase, \(q=12\) gives the
exact identity

\[
\|12(\sqrt2-1)\|_{\mathbb Z}
=17-12\sqrt2=\frac{1}{17+12\sqrt2}.
\tag{6.15}
\]

Here the nearest integer is \(5\). This illustrates a small arithmetic
phase mismatch without relying on a decimal cosine value. It is
\textbf{not} presented as a completed run of the rational-input
certificate at the real target (6.2).

The contrasting regime is Liouville approximation. Recall that an
irrational \(\alpha\) is a \emph{Liouville number} if, for every
positive integer \(M\), there are infinitely many rational
approximations \(a/q\) with \(q\to\infty\) and
\(|\alpha-a/q|<q^{-(M+1)}\). At those denominators,

\[
\|q\alpha\|_{\mathbb Z}<q^{-M},
\qquad
\Delta_q<2\pi|\beta|q^{-(M+1)}.
\tag{6.16}
\]

For each fixed \(M\), Corollary 6.3 therefore gives, along an unbounded
sequence of denominators,

\[
p_q\ge(M+1)\log_2q-O_\beta(1).
\tag{6.17}
\]

Because \(M\) is arbitrary, \(\limsup_{q\to\infty}p_q/\log_2q=+\infty\).
Equivalently, the inverse norm in (6.9) grows faster than any prescribed
fixed power of \(q\) along suitable subsequences. This does not mean
that the resolvent grows monotonically with matrix size, or that every
large denominator is resonant. The conclusion is explicitly
\emph{subsequential}.

The familiar Liouville series \(\sum_{j\ge1}10^{-j!}\) provides an exact
mathematical example of such a phase in \((0,1)\). Its defining series
and rational partial sums do not, by themselves, furnish a validated
rational enclosure of the cosine target required by our executable
certificate.

\subsection{6.4. What this example establishes---and what it does
not}\label{what-this-example-establishesand-what-it-does-not}

The constant-chain analysis concerns a real phase target
\(\tau_\alpha\), whereas the executable producer of Section 4 takes a
rational target. A computational instance with an irrationally specified
phase therefore needs another, independently certified input: for
example, rational \(r\) and \(\rho\ge0\) for which

\[
|\tau_\alpha-r|\le\rho.
\tag{6.18}
\]

If the coefficients \(c,\beta\) are rational and the rational producer
certifies \(L\le\Delta(T_{q-1},r)\le U\), then Theorem 5.1 transports
the conclusion to the exact phase target:

\[
\max\{0,L-\rho\}
\le\Delta_q\le U+\rho.
\tag{6.19}
\]

The claim (6.18) must be earned by validated bounds on the phase,
\(\pi\), and cosine, or by another rigorous argument. An unqualified
decimal approximation to \(\tau_\alpha\) is not such a certificate.
Likewise, the cost of explicitly forming and processing a very large
constant chain is separate from the closed-form phase identity (6.3).

Three claims are consequently kept distinct. The identity (6.3) and
comparison (6.7) are classical analytic properties of an exactly
solvable chain. Proposition 6.2 is a sharp statement for the information
model (6.10). The producer of Section 4 is a different exact-rational
computation, whose soundness follows from the independent Sturm checker.
The arithmetic example explains why small certified margins matter, but
does not establish a universal lower bound on computational complexity,
a verified floating-point eigensolver, or a completed irrational-target
front end.

This distinction returns us to the central question of the paper. At a
resonant denominator the spectral distance may be extremely small even
though the matrix coefficients are simple. What matters for a claim of
safe invertibility is not that a numerical eigenvalue seems separated
from the target, but that the proposed separation comes with evidence
which a second computation can check.

\section{7. Relation to existing certification methods and the trust
boundary}\label{relation-to-existing-certification-methods-and-the-trust-boundary}

\subsection{7.1. Why a target-specific
witness?}\label{why-a-target-specific-witness}

A finite Jacobi characteristic polynomial could be submitted to a
general exact real-root isolator, its roots enclosed, and the nearest
enclosing interval compared with the desired energy. The resulting
method can be fully rigorous. There is no claim here that an exact
target-distance certificate computes information unattainable from those
general algorithms. Our distinction is one of \emph{question, evidence,
and interface}: we do not seek an ordered list of all eigenvalue
enclosures, nor do we need an index for the closest root. The checker
requires only an eigenvalue-free target neighborhood and one occupied
neighborhood, expressed as rational Sturm-count conditions. The output
is the quantitative margin that the resolvent and residual inequalities
consume directly.

The tridiagonal structure supplies the characteristic continuant by a
three-term recurrence. For a fixed rational test point, its values
require a finite sequential computation involving rational additions and
multiplications. The source construction uses the same recurrence to
obtain a positive root-free seed and a finite occupied outer interval.
None of these observations entails that rational bit lengths remain
small or that the present producer is competitive with mature numerical
eigensolvers. The distinction between arithmetic-operation count, exact
integer growth, and elapsed time is maintained throughout.

Fernando {[}1{]} develops guaranteed eigenvalue inclusions for symmetric
tridiagonal matrices from an exact-inertia calculation carried out with
suitable directed-rounding floating-point arithmetic. That is an
important precedent for rigorous tridiagonal spectral computation. Our
exact-rational setting has a different trust boundary: the checker
recomputes rational signs and count comparisons and invokes a formal
theorem tying those counts to the spectrum. We do not certify a
floating-point inertia calculation or claim that rational arithmetic is
preferable for the same engineering workload.

Li, Passmore and Paulson {[}2{]} exemplify the separation of certificate
production from formal verification for univariate polynomial problems
in Isabelle/HOL. The Lean \texttt{hex-real-roots} family {[}3{]} now
goes further in another direction: certified general integer-polynomial
root isolation, half-open Sturm witnesses, completeness and an
explicitly designed faster search strategy. The present paper does not
claim novelty in that architectural separation, in zero-skipping sign
variations, or in half-open interval conventions. Its proposed value is
the \emph{composable specialized object}: a finite four-endpoint witness
for a specified Jacobi spectral margin, a constructive exact-data path
to such a witness, and a direct link to inverse norms and represented
perturbation budgets.

This specialization is useful only if its execution costs and
source-level guarantees survive inspection. In particular, one should
not confuse a total program with a practical one. Repeated bisection has
width \((u_0-\ell_0)/2^k\); for a requested absolute width
\(\varepsilon\), the natural refinement count is the least integer
\(k\ge0\) satisfying

\[
2^k\varepsilon\ge u_0-\ell_0,
\qquad
k=\max\left\{0,\left\lceil
\log_2\frac{u_0-\ell_0}{\varepsilon}\right\rceil\right\}.
\tag{7.1}
\]

This is an iteration-count identity, \textbf{not} a bit-complexity
theorem. The executable stopping rule can use integer ceiling and
base-two logarithm routines to avoid floating-point evaluation of (7.1).
The implementation must separately show that the chosen count satisfies
the inequality and that rational costs remain manageable. A much larger
terminating bound remains mathematically correct but should not be
advertised as an efficient implementation.

\subsection{7.2. What an accepted artifact
establishes}\label{what-an-accepted-artifact-establishes}

For a matrix encoded by finite rational diagonal and off-diagonal lists,
a rational target \(\tau\), positive rational tolerance \(\varepsilon\),
and a returned four-endpoint record \(W\), the verifier's mathematical
contract is

\[
\begin{split}
&\operatorname{Check}(J,\tau,W)=\mathrm{true},\\
&0<L(W,\tau)\le U(W,\tau),\\
&U(W,\tau)-L(W,\tau)\le\varepsilon.
\end{split}
\tag{7.2}
\]

These are checks on data, not claims about which search algorithm
generated the data. Together with the counting theorem they establish
the spectral-distance enclosure of Theorem 3.1. The full implementation
should expose the results in a stable, finite-input format and make it
possible to rerun the checker on the exact serialized output. A warm
build of individual theorem declarations is not equivalent to such a
released, end-to-end evidence object. Conversely, the mathematical
theorem remains valid even if a particular producer is inefficient or
replaced by a better search strategy.

The trust boundary also does not absorb every external numerical
statement. A residual estimate must be validated by its own computation;
a bound on perturbation of the input coefficients must be supplied or
checked; and an irrationally specified target requires a certified
rational enclosure before the rational procedure can be applied to it.
Section 5 states the analytic transport theorem with those hypotheses
visible. Section 6's arithmetic precision threshold belongs to a
specified reporting model; it does not impose a universal precision
lower bound on other certificate formats.

\subsection{7.3. Scope and limitations}\label{scope-and-limitations}

The restriction to unreduced rational Jacobi matrices is substantive. A
reduced chain can be decomposed into unreduced blocks, but such
preprocessing must not silently be inferred from a checker that rejects
zero couplings. The diagonal and off-diagonal prefixes determine the
physical matrix; output must not depend on arbitrary coefficients at
unused indices of an internal extension. The singular case is reported
separately from a positive spectral margin; a totalized matrix-inverse
operation is not an appropriate substitute for that distinction.

The paper's rigorous perturbation and inverse-norm consequences are
classical matrix theory, assembled to make the certificate useful. They
are not offered as new Weyl-type inequalities. Its fixed-phase
arithmetic examples reveal a demanding spectral scale, but their
analytic asymptotics are not proof of the certificate's runtime or of a
general model of numerical information. Such separations are essential
to judging both the contribution and its limitations accurately.

\section{References}\label{references}

{[}1{]} K. V. Fernando, ``Computation of exact inertia and inclusions of
eigenvalues (singular values) of tridiagonal (bidiagonal) matrices,''
\emph{Linear Algebra and its Applications} \textbf{422} (2007), 77--99.
https://doi.org/10.1016/j.laa.2006.09.008.

{[}2{]} W. Li, G. O. Passmore, and L. C. Paulson, ``Deciding univariate
polynomial problems using untrusted certificates in Isabelle/HOL,''
\emph{Journal of Automated Reasoning} \textbf{62} (2019), 69--91.
https://doi.org/10.1007/s10817-017-9424-6.

{[}3{]} Lean development team, \emph{hex-real-roots: certified real-root
isolation for dense integer polynomials}, open-source specification and
Mathlib companion, https://github.com/leanprover/hex-real-roots and
https://github.com/leanprover/hex-real-roots-mathlib (consulted October
2026; release commit to be pinned with the final companion).

{[}4{]} O. A. Marques, E. J. Riedy, and C. Vömel, ``Benefits of IEEE-754
features in modern symmetric tridiagonal eigensolvers,'' \emph{SIAM
Journal on Scientific Computing} \textbf{28} (2006), 1613--1633.
https://doi.org/10.1137/050641624.

\end{document}
